% ID: SEL-D0016-BB-E03
% ESTADO: validada. S017. Revisión matemática, didáctica y visual.
\textbf{a) Derivadas.}\par\nopagebreak[4]
Aplicamos la regla de la cadena al logaritmo y la del cociente al término racional:
\[
f'(x)=\frac{6x}{3x^2-3}+
\frac{-2(x+2)-(1-2x)}{(x+2)^2}
=\boxed{\frac{2x}{x^2-1}-\frac5{(x+2)^2}}.
\]
El dominio exige $|x|>1$ y $x\ne-2$. Para $g$, usamos cadena en la exponencial y producto en el segundo sumando:
\[
\boxed{g'(x)=6x^2e^{x^3}+2x(3x+4)^3+9x^2(3x+4)^2.}
\]
Esta derivada existe en todo $\mathbb R$.

\textbf{b) Tangentes y corte.}\par\nopagebreak[4]
Para $h$, $h(1)=2$ y $h'(1)=2$. Para $p$, $p(1)=0$ y $p'(x)=2/(x+1)^2$, de donde $p'(1)=1/2$. Por la ecuación punto-pendiente, las tangentes son
\[
\boxed{y=2+2(x-1)=2x},\qquad \boxed{y=\tfrac12(x-1)}.
\]
Igualando $2x=(x-1)/2$ se obtiene $3x=-1$. El corte es $\boxed{(-1/3,-2/3)}$.
